AI agents increasingly act on users' behalf through tools supplied by third-party Model Context Protocol (MCP) servers, gaining access to email, files, shells, and credentials. Users and clients approve these tools once, based on their names and descriptions. Yet servers keep changing after approval. Prior work shows that clients do not re-request approval when a tool changes, but all existing MCP measurements analyze servers at a single point in time, so it is unknown how often approved tools actually change, and whether anyone can see it.

We present the first longitudinal study of post-approval drift in MCP servers. We analyze 35,725 sampled versions from the release histories of 2,571 open-source servers, extract each tool's description and code-level capabilities at every version, and validate our measurements against a labeled sample. About two-thirds of servers change their tool descriptions after release, while capabilities such as shell execution, file writes, and secret access are frequently added with no change to the description users approved. Most changes are benign, which is itself a problem: risky changes are buried in constant routine updates. Attacks studied in prior work, such as hidden Unicode instructions, are virtually absent in practice; instead, we find servers that delegate risky actions to the AI agent itself, outside the view of code scanners.

Simulating client policies on real histories, we show that re-approving on every update would prompt users many times per server per year, while re-approving on description changes misses every silent capability gain. We argue that approval must be tied to code-derived capabilities, and we derive concrete recommendations for MCP clients and the protocol.
